\chapter{Introduction}
\label{ch:into}

This chapter will explain the context of our project, the multi-agent challenges that we are tackling,
the engineering objectives we will pursue, the methodology and solution approaches we will use,
and the project management timeline for the project.

%%%%%%%%%%%%%%%%%%%%%%%%%%%%%%%%%%%%%%%%%%%%%%%%%%%%%%%%%%%%%%%%%%%%%%%%%%%%%%%%%%%
\section{Project Presentation}
\label{sec:proj_pres}

\subsection{Background}
\label{sec:into_back}

The concept of Deep Reinforcement Learning (DRL) has brought a new level of innovation to 
autonomous agent development. Instead of creating manual rules for every situation, we allow 
the agent to learn through its own experiences in a given environment. Rather than relying 
on a programmer to predict every single possible scenario and create an appropriate rule for 
each scenario, the agent develops effective strategies through experimentation with the help 
of a reward signal that reflects the quality of the decisions made. Successful examples of 
DRL include AlphaGo~\citep{silver2016mastering} and deep Q learning networks which played 
games like Atari~\citep{mnih2015humanlevel} where both used reinforcement learning to achieve 
strategic intuition and behavioral complexity that could never be produced by a human programmer 
through hard coding. These developments led to a great deal of interest in extending DRL to 
even more challenging domains such as robot locomotion and real-time strategy gaming.

MARL takes this idea even further than DRL by introducing multiple agents into one environment 
learning simultaneously. This reality adds an entirely new level of challenge to learning; 
now not only must an agent learn how to work by itself, but also learn how to change its 
behavior based on the actions of the other agents in the environment. Depending on the type 
of tasks being tested, it is possible for a number of different types of structures to exist 
such as;

%%%%%%%%%%%%%%%%%%%%%%%%%%%%%%%%%%%%%%%%%%%%%%%%%%%%%%%%%%%%%%%%%%%%%%%%%%%%%%%%%%%
\subsection{Problem Statement}
\label{sec:intro_prob_art}

Our main goal is to explore how we can use only reinforcement learning in training multiple 
autonomous agents to work, cooperate and compete in a physics-based environment that has 
continuous action space. There are four specific behaviors that four physically based agents 
need to learn from nothing:

\begin{enumerate}
    \item The ability to control their joints and travel around in the continuous three 
    dimensional pitch.
    \item Locate, chase and intercept a moving soccer ball.
    \item Coordinate with a teammate to carry the ball toward the goal.
    \item Defend their own goal against two other opponents and attempt to score against them.
\end{enumerate}

Due to this nature of the task, it is very difficult to accomplish. Goals will not happen 
often when executing random policies, resulting in a sparse reward signal. Moreover, state 
space contains continuous position, direction and velocity, and partial-timed learning will 
violate the Markov property for each individual agent.

%%%%%%%%%%%%%%%%%%%%%%%%%%%%%%%%%%%%%%%%%%%%%%%%%%%%%%%%%%%%%%%%%%%%%%%%%%%%%%%%%%%
\subsection{Aims and Objectives}
\label{sec:intro_aims_obj}

Our goal is to investigate the design, training, and evaluation of a cooperative 2-on-2 (2v2) 
soccer team using a continuous three-dimensional (3D) MuJoCo simulation with different analysis 
methods and train the agents using Proximal Policy Optimization (PPO), reward shaping wrappers, 
and spatial analysis.

\textbf{Objectives:}

\begin{enumerate}
    \item Create and set up the MuJoCo 2v2 soccer environment by utilizing the PettingZoo 
    and Shimmy wrappers so that each player in the 4-player simulation can interact in a 
    stable manner across all four players in the physical environment.

    \item Design a custom reward shaping wrapper to change goal-based (i.e., sparse) rewards 
    to be dense (i.e., immediate) based on the distance and/or velocity of the agent to the 
    goal so that an agent can receive useful and meaningful feedback during an episode, instead 
    of just when the agent scores.

    \item Use PPO along with a shared policy network and/or parameter sharing across all agents 
    in order to train them to be able to exhibit cooperative behavior consistently and reduce 
    sample complexity to improve the stability of the training in a non-stationary multi-agent 
    environment.

    \item Create a soccer analytics pipeline to analyze soccer data by extracting and 
    visualizing the following metrics: team possession \%, spatial spread index, and hexbin 
    density heatmap, to provide a quantitative basis for evaluating path developed tactical 
    organization and positional role separation among agents.
\end{enumerate}

%%%%%%%%%%%%%%%%%%%%%%%%%%%%%%%%%%%%%%%%%%%%%%%%%%%%%%%%%%%%%%%%%%%%%%%%%%%%%%%%%%%
\subsection{Solution Approach}
\label{sec:intro_sol}

Four major components make up our engineering pipeline:

\begin{itemize}
    \item \textbf{Stage 1 --- Environment Integration:} To create a stable and fast parallel 
    vectorized interface with multiple instances of the MuJoCo (2v2) soccer environment, we 
    use PettingZoo and Shimmy compatibility wrappers. This not only gives us a single interface 
    to all environments, but it also allows us to collect experience batches as quickly and 
    efficiently as possible.

    \item \textbf{Stage 2 --- Reward Wrapper Logic:} We create a new class 
    (\texttt{SoccerRewardShapingWrapper}) to enhance the sparse goal-based reward signal with 
    dense, physics-based feedback about how far away an agent is from the ball and how fast 
    it is moving toward the goal, making sure that all agents receive meaningful feedback 
    every step of the way.

    \item \textbf{Stage 3 --- Policy Optimization:} We use a single shared PPO (actor-critic) 
    network for all teammates and use Google Colab to train using 1 T4 GPU and multiprocessing 
    to speed up the collection of experiences and updates to the policy across all the 
    environments.

    \item \textbf{Stage 4 --- Spatial and Game Analysis:} After training is finished, we run 
    dedicated post-training evaluation episodes to log the position of the agent during each 
    game. This positional data is then processed to create 2D hexbin density heatmaps showing 
    coverage of the pitch and possession rate calculations for each team to quantify their 
    ball control relatively.
\end{itemize}

% ================================================================
\section{Project Management}
\label{sec:proj_mgmt}
% ================================================================

We scheduled our project over a three-week period. Table~\ref{tab:gantt} outlines the key phases and timelines of our work.

\begin{table}[!ht]
    \centering
    \caption{Project management timeline}
    \label{tab:gantt}
    \begin{tabular}{llc}
        \toprule
        \textbf{Phase} & \textbf{Tasks} & \textbf{Duration} \\
        \midrule
        1. Setup         & Dependency installation, MuJoCo configuration, API wrapper checks & Week 1 \\
        2. Modeling      & Reward shaping design, state-action interface analysis & Week 1--2 \\
        3. Implementation & Environment wrappers, training scripts, Colab notebook structure & Week 2 \\
        4. Training      & PPO model training on Colab, checkpoint saving & Week 2--3 \\
        5. Analysis      & Team possession metrics, spatial coverage heatmaps & Week 3 \\
        6. Compilation   & Writing the LaTeX report and finalizing code comments & Week 3 \\
        \bottomrule
    \end{tabular}
\end{table}

% ================================================================
\section{Summary}
\label{sec:intro_summary}
% ================================================================

This chapter describes the framework, problem statement and goals for our multi-agent soccer 
project. We established four clearly defined objectives: integrating into an environment, 
shaping the reward function, optimising the policy, and performing spatial analysis. We also 
produced a complete solution workflow and timeline. These objectives demonstrated the complexity 
of the project due to the effects of continuous action spaces, non-stationary multi-agents, 
and sparse reward signals all typical of environments with goal scoring. In Chapter~\ref{ch:lit_rev}, 
we will provide theoretical basis and literature review associated with reinforcement learning, 
such as Markov Decision Processes, and value-based methods of reinforcement learning; as well 
as algorithms of continuous actor-critic methods that serve as a foundation for our PPO based 
training.